% Part: first-order-logic 
% Chapter: axiomatic-deduction 
% Section: provability

% verification of properties of provability needed for maximally 
% consistent sets in the completeness chapter.

\documentclass[../../../include/open-logic-section]{subfiles}

\begin{document}

\iftag{FOL}
      {\olfileid{fol}{axd}{prv}}
      {\olfileid{pl}{axd}{prv}}

\olsection{\usetoken{S}{derivability}ના ગુણધર્મો}

\begin{prop}[એકદિશવર્ધિતા] જો $\Gamma \subseteq \Delta$ અને $\Gamma \Proves !A$
હોય, તો $\Delta \Proves !A$ પણ હોય. \end{prop}

\begin{proof} કોઈ પણ સાન્ત $\Gamma_0 \subseteq \Gamma$ એ
$\Delta$નો પણ સાન્ત ઉપગણ છે. તેથી
$\Gamma_0 \Proves !A$ની !!a{derivation}
$\Delta \Proves !A$ પણ બતાવે છે. \end{proof}

\begin{prop}\ollabel{prop:provability} $\Gamma \Proves !A$ જો અને માત્ર જો
$\Gamma\cup \{\lnot!A \}$ અસુસંગત હોય. \end{prop}

\begin{prop} \begin{enumerate}
\item \ollabel{prop:provability-contr} જો $\Gamma \Proves !A$ અને $\Gamma
\cup \{ !A\} \Proves \lfalse$, તો $\Gamma$ અસુસંગત છે.

\item \ollabel{prop:provability-lnot} જો $\Gamma \cup \{!A\}
\Proves \lfalse$, તો $\Gamma \Proves \lnot !A$.

\item \ollabel{prop:provability-exhaustive} જો $\Gamma \cup \{!A\}
\Proves \lfalse$ અને $\Gamma \cup \{\lnot !A\} \Proves
\lfalse$, તો $\Gamma \Proves \lfalse$.

\tagitem{prvOr}{\ollabel{prop:provability-lor-left} જો $\Gamma \cup \{!A\}
\Proves \lfalse$ અને $\Gamma \cup \{!B\} \Proves
\lfalse$, તો $\Gamma \cup \{!A \lor !B\} \Proves \lfalse$.}{}

\tagitem{prvOr}{\ollabel{prop:provability-lor-right} જો $\Gamma
\Proves !A$ અથવા $\Gamma \Proves !B$, તો $\Gamma
\Proves !A \lor !B$.}{}

\tagitem{prvAnd}{\ollabel{prop:provability-land-left} જો $\Gamma
\Proves !A \land !B$, તો $\Gamma \Proves !A$ અને
$\Gamma \Proves !B$.}{}

\tagitem{prvAnd}{\ollabel{prop:provability-land-right} જો $\Gamma
\Proves !A$ અને $\Gamma \Proves !B$, તો $\Gamma
\Proves !A \land !B$.}{}

\tagitem{prvIf}{\ollabel{prop:provability-mp} જો $\Gamma \Proves
!A$ અને $\Gamma \Proves !A \lif !B$, તો $\Gamma
\Proves !B$.}{}

\tagitem{prvIf}{\ollabel{prop:provability-lif} જો $\Gamma \Proves
\lnot !A$ અથવા $\Gamma \Proves !B$, તો $\Gamma \Proves
!A \lif !B$.}{} \end{enumerate} \end{prop}

\begin{proof}
\begin{enumerate}
\item માનો કે $\Gamma_0 \cup \{!A\} \Proves \lfalse$.

\begin{derivation} 
1. & $\Gamma_0 \cup \{!A\} \Proves \lfalse$ & ધારણા \\
2. & $\Gamma_1 \Proves !A$ & ધારણા\\
3. & $\Gamma_0 \cup \Gamma_1 \Proves \lfalse$ & કટ \\
\end{derivation}

કારણ કે $\Gamma_0 \subseteq \Gamma$ અને $\Gamma_1 \subseteq \Gamma$,
તેથી $\Gamma_0 \cup \Gamma_1 \subseteq \Gamma$; આથી $\Gamma \Proves \lfalse$.

\item માનો કે $\Gamma_0 \cup\{!A\} \Proves \lfalse$.

\begin{derivation} 
1. & $\Gamma_0 \cup\{!A\} \Proves \lfalse$ & ધારણા \\
2. & $\Gamma_0 \Proves !A \lif \lfalse$ & નિગમન પ્રમેય\\
3. & $\Gamma_0 \Proves (!A \lif \lfalse) \lif \lnot !A$ & Ax0\\
4. & $\Gamma_0 \Proves \lnot !A$ & MP 2, 3 \\
\end{derivation}
\sourcecorrection{OLMD-174}{મૂળના આ બીજા
કિસ્સામાં વિધાન
$\Gamma\cup\{!A\}\Proves\lfalse$
છે, પરંતુ સાબિતી
$\Gamma_0\cup\{\lnot !A\}\Proves\lfalse$
થી શરૂ થઈ
$!A$ પર પહોંચે છે.
અહીં ધારણા અને
પછીના ચાર પગલાં
વિધાન મુજબ
$\lnot !A$ કાઢવા
સુધાર્યાં છે.}

\item માનો કે $\Gamma_0 \cup \{ !A \} \Proves \lfalse$ અને
$\Gamma_1 \cup \{\lnot !A \} \Proves \lfalse$.

\begin{derivation} 
1. & $\Gamma_0 \cup\{!A\} \Proves \lfalse$ & ધારણા \\
2. & $\Gamma_0 \Proves !A \lif \lfalse$ & નિગમન પ્રમેય \\
3. & $\Gamma_1 \cup \{\lnot !A \} \Proves \lfalse$ & ધારણા\\
4. & $\Gamma_0 \Proves (!A \lif \lfalse) \lif \lnot !A$ & Ax0\\
5. & $\Gamma_0 \Proves \lnot !A$ & MP 2, 4\\
6. & $\Gamma_0 \cup \Gamma_1 \Proves \lfalse$ & કટ 3, 5 \\
\end{derivation}

કારણ કે $\Gamma_0 \subseteq \Gamma$ અને $\Gamma_1 \subseteq \Gamma$,
તેથી $\Gamma_0 \cup \Gamma_1 \subseteq \Gamma$.
આથી $\Gamma \Proves \lfalse$.
\sourcecorrection{OLMD-175}{મૂળમાં
આ ત્રીજા કિસ્સાની
પ્રારંભિક સાન્ત
ધારણા માટે
$\Gamma$ લખાયું છે,
પણ પછીની નિષ્પત્તિ
$\Gamma_0$ વાપરે છે.
સાન્ત સાક્ષી
$\Gamma_0\subseteq\Gamma$
સ્પષ્ટ કરવા પ્રારંભિક
સૂત્ર સુધાર્યું છે.}

% prop:provability-lor-left 
\tagitem{defOr}{}{
\iftag{probOr}{કસરત.}{કસરત.}}

% prop:provability-lor-right 
\tagitem{defOr}{}{ \iftag{probOr}{કસરત.}{
માનો કે $\Gamma_0 \Proves !A$.

\begin{derivation} 
1. & $\Gamma_0 \Proves !A$ & ધારણા \\
2. & $\Gamma_0 \Proves !A \lif (!A \lor !B)$ & Ax0\\ 
3. & $\Gamma_0 \Proves !A \lor !B$ & MP 1, 2 \\
\end{derivation}

તેથી $\Gamma \Proves !A \lor !B$.
$\Gamma \Proves !B$ હોય ત્યારે
સાબિતી સમાન છે.}}

% prop:provability-land-left 
\tagitem{defAnd}{}{ \iftag{probAnd}{કસરત}{
માનો કે $\Gamma_0 \Proves !A \land !B$

\begin{derivation} 
1. & $\Gamma_0 \Proves !A \land !B$ & ધારણા \\
2. & $\Gamma_0 \Proves (!A\land !B) \lif !A $ & Ax0\\ 
3. & $\Gamma_0 \Proves !A$ & MP 1, 2 \\
\end{derivation}
\sourcecorrection{OLMD-176}{મૂળની
ત્રીજી પંક્તિમાં
પૂર્વધારણા અને
બીજી પંક્તિના
પ્રેરણમાંથી
$!A\lor !B$
કાઢ્યું છે. મોડસ
પોનેન્સથી અહીં
$!A$ મળે છે;
તે મુજબ સૂત્ર
સુધાર્યું છે.}

આથી $\Gamma \Proves !A$.
તેવી જ !!{derivation}
$\Gamma \Proves !B$
પણ બતાવે છે.}}

% prop:provability-land-right
\tagitem{defAnd}{}{\iftag{probAnd}{કસરત.}{કસરત.}}

% prop:provability-mp 
\tagitem{defIf}{}{ \iftag{probIf}{કસરત.}{ માનો કે
$\Gamma_0 \Proves !A$ અને $\Gamma_0 \Proves !A \lif !B $ બંને છે.

\begin{derivation} 
1. & $\Gamma_0 \Proves !A$ & ધારણા \\
2. & $\Gamma_0 \Proves !A \lif !B $ & ધારણા\\
3. & $\Gamma_0 \Proves !B$ & MP 1, 2\\
\end{derivation}

કારણ કે $\Gamma_0 \cup \Gamma_1 \subseteq \Gamma$,
આથી $\Gamma \Proves !B$.
\sourcecorrection{OLMD-177}{મૂળના
આ કિસ્સામાં
$\Gamma_1$ ક્યાંય
નક્કી કરાયું નથી.
નિષ્કર્ષ માટે
પહેલેથી પસંદ
કરાયેલ સાન્ત
$\Gamma_0\subseteq\Gamma$
જ પૂરતું છે;
વધારાના
$\Gamma_1$ની
જરૂર નથી.}}}

% prop:provability-lif 
\tagitem{defIf}{}{\iftag{probIf}{કસરત.}{પહેલાં
માનો કે $\Gamma \Proves \lnot !A$.
    
\begin{derivation} 
1. & $\Gamma_0 \Proves \lnot !A$ & ધારણા \\
2. & $\Gamma_0 \Proves \lnot !A \lif (!A \lif !B) $ & Ax0\\ 
3. & $\Gamma_0 \Proves !A \lif !B$ & MP 1,
2\\ \end{derivation}

$\Gamma \Proves !B$ માટેની સાબિતી સમાન છે:

\begin{derivation} 
1. & $\Gamma_0 \Proves !B$ & ધારણા \\
2. & $\Gamma_0 \Proves !B \lif (!A \lif !B) $ & Ax0\\ 
3. & $\Gamma_0 \Proves !A \lif !B$ & MP 1, 2\\
\end{derivation}}}

\end{enumerate} 
\end{proof}

\begin{probtag}{probOr,probAnd,probIf} આની સાબિતી પૂરી કરો:
\olref[fol][axd][prv]{prop:provability}. 
\end{probtag}

\begin{thm}[સર્વત્રીકરણ] \ollabel{thm:Generalization} જો $\Gamma \Proves
!A$ અને $x$ એ $\Gamma$નાં કોઈ પણ !!{formula}માં મુક્ત ન હોય,
તો $\Gamma \Proves \lforall[x][!A]$.
\end{thm}

\begin{proof} $\mathsf{Thm}_\Gamma$ પરના આગમનથી:
જો $!A$ સ્વયંસિદ્ધિ હોય, તો
$\lforall[x][!A]$ પણ હોય.
જો $!A\in \Gamma$, તો $x$ એ $!A$માં મુક્ત
નથી; તેથી $!A \lif \lforall[x][!A]$
સ્વયંસિદ્ધિ છે (\textbf{Ax3}) અને MPથી
$\Gamma \Proves \lforall[x][!A]$ મળે.
હવે માનો કે $\Gamma \Proves !B$ અને
$\Gamma \Proves !B \lif !A$ હોવાથી
MPથી $!A$ મળે છે. આગમન-પરિકલ્પનાથી
$\Gamma \Proves \lforall[x][!B]$ અને
$\Gamma \Proves \lforall[x][( !B \lif !A)]$.
પરંતુ \textbf{Ax2}થી આ પણ મળે:
\[ \Gamma \Proves
\lforall[x][( !B \lif !A)] \lif (\lforall[x][ !B] \lif \lforall[x][!A]), \]
અને MPના બે ઉપયોગથી ઇચ્છિત
$\Gamma \Proves \lforall[x][!A]$ મળે. \end{proof}

\begin{thm}\ollabel{thm:WeakGen} (\emph{અચળો પર નબળું સર્વત્રીકરણ})
જો $\Gamma \Proves !A$ અને $c$ એ $\Gamma$માં
ન આવતું !!a{constant} હોય, તો $!A$માં
ન આવતું કોઈ ચલ $x$ એવું મળે કે
$\Gamma \Proves \lforall[x][\Subst{!A}{x}{c}]$;
આ સાબિતીમાં $c$ આવતું નથી.
\end{thm}

\begin{proof} $!A_1,\ldots,!A_n$ એ $\Gamma$માંથી
$!A$ની સાબિતી માનો, જેથી $!A=!A_n$.
$!A_1,\ldots,!A_n$માં ન આવતું
ચલ $x$ પસંદ કરો અને
નવો અનુક્રમ લો:
$\Subst{!A_1}{x}{c},\ldots,\Subst{!A_n}{x}{c}.$
આવો અનુક્રમ $\Gamma$માંથી
$\Subst{!A}{x}{c}$ની સાબિતી છે.
ખરેખર, દરેક $i=1,\ldots,n$ માટે:
\begin{itemize}
\item જો $!A_i$ સ્વયંસિદ્ધિ હોય, તો
$\Subst{!A_i}{x}{c}$ પણ સ્વયંસિદ્ધિ છે;
\item જો $!A_i \in \Gamma$, તો
$c$ એ $\Gamma$માં નથી,
એટલે $\Subst{!A_i}{x}{c} = !A_i \in \Gamma$;
\item જો $!A_i$ એ $!A_j$
અને $!A_j \lif !A_i$માંથી મળે,
તો $\Subst{!A_j}{x}{c}$ અને
$\Subst{(!A_j \lif !A_i)}{x}{c}
= \Subst{!A_j}{x}{c} \lif \Subst{!A_i}{x}{c}$
પર MP લગાવી $\Subst{!A_i}{x}{c}$ મળે.
\end{itemize}
\sourcecorrection{OLMD-178}{મૂળની
સાબિતીમાં દરેક
$i$ માટે કિસ્સા
જુદાં પાડવામાં
આવે છે, છતાં
પહેલી શરતમાં
$!A$ લખાયું છે.
તેને સંબંધિત
પંક્તિ $!A_i$
કરી છે.}
નવા અનુક્રમમાં !!{constant} $c$
હવે આવતું નથી તે સ્પષ્ટ છે.
હવે $\Gamma$નાં જે !!{formula}s
$\Subst{!A}{x}{c}$ની !!{derivation}માં
આવે તેમનો ગણ $\Gamma'$ લો.
ત્યારે $x$ એ $\Gamma'$નાં
કોઈ પણ !!{formula}માં મુક્ત નથી
અને $\Gamma' \Proves \Subst{!A}{x}{c}$.
સર્વત્રીકરણથી
(\olref{thm:Generalization})
$\Gamma' \Proves \lforall[x][\Subst{!A}{x}{c}]$,
અને એકદિશવર્ધિતાથી ઇચ્છિત
$\Gamma \Proves \lforall[x][\Subst{!A}{x}{c}]$ મળે.
\end{proof}

\begin{explain}
આપણો હેતુ \olref{thm:WeakGen}ની
$x$ એ $!A$માં બિલકુલ ન આવતું
હોવાની શરતને બે નબળી શરતોથી
બદલવાનો છે: $x$ એ $!A$માં
મુક્ત ન હોય અને $!A$માં
$c$ની જગ્યાએ મૂકવા !!{free for}
હોય. બંધ !!{variable} બદલીને
આ કરી શકાય છે; તેથી પહેલાં
એ જ સાબિત કરીશું.
\end{explain}

\begin{lem}
\ollabel{lem:free-for} 
જો $x$ એ $!A$માં $c$ની જગ્યાએ
મૂકવા મુક્ત હોય અને $y$ એ $!A$માં
મુક્ત ન હોય, તો $x$ એ
$\Subst{!A}{y}{c}$માં $y$ની
જગ્યાએ મૂકવા !!{free for} છે.
\end{lem}

\begin{proof} જો $x$ એ
$\Subst{!A}{y}{c}$માં
$y$ની જગ્યાએ મૂકવા
!!{free for} ન હોય,
તો $\Subst{!A}{y}{c}$માં
$y$નું કોઈ મુક્ત આગમન
$\lforall[x]$ પરિમાણકના
વ્યાપમાં પડે. પરંતુ
પરિકલ્પના મુજબ $y$
એ $!A$માં મુક્ત નથી;
તેથી આવા બધા આગમનો
$\Subst{}{y}{c}$ અવેજીથી
આવે છે. એટલે $c$નું
કોઈ આગમન $\lforall[x]$ના
વ્યાપમાં પડે અને $x$
એ $!A$માં $c$ની
જગ્યાએ મૂકવા
!!{free for} ન રહે.
\sourcecorrection{OLMD-179}{મૂળની
સાબિતીના પ્રથમ વાક્યમાં
$\Subst{!A}{y}{x}$
લખાયું છે, પરંતુ
પ્રમેયનું વિધાન અને
આગળનાં વાક્યો
$\Subst{!A}{y}{c}$
વિશે છે. એ જ
અવેજી પ્રથમ
વાક્યમાં સુધારી છે.}
\end{proof}

\begin{lem}[બંધ ચલનું પરિવર્તન]
\ollabel{lem:ChangeBdVar} 
જો $x$ અને $y$ એ $!A$માં મુક્ત
ન હોય અને બંને $!A$માં $c$ની
જગ્યાએ મૂકવા !!{free for}
હોય, તો $\Proves \lforall[x][\Subst{!A}{x}{c}] \equiv
\lforall[y][\Subst{!A}{y}{c}]$;
અને સાબિતીમાં $c$ આવતું નથી.
\end{lem}

\begin{proof}
પરિકલ્પનાથી $y$ એ $!A$માં
$c$ની જગ્યાએ મૂકવા !!{free for}
છે અને $x$ એ $!A$માં મુક્ત
નથી; તેથી અગાઉના
\olref{lem:free-for}થી $y$ એ
$\Subst{!A}{x}{c}$માં
$x$ની જગ્યાએ મૂકવા
!!{free for} છે. આથી
$\lforall[x][\Subst{!A}{x}{c}] \lif \Subst{!A}{x}{c} \Subst{}{y}{x}$
સ્વયંસિદ્ધિ છે (\textbf{Ax1}).
\sourcecorrection{OLMD-180}{મૂળમાં
સર્વવાચી પરિમાણક
આખા પ્રેરણ
પર મૂકાયો છે,
પરંતુ \textbf{Ax1}
માટે જરૂરી સ્વરૂપ
સર્વત્રીકૃત પૂર્વાંગ
પછી તેના દૃષ્ટાંત
તરફના પ્રેરણનું છે.
અહીં કૌંસ અને
પ્રેરણનું સ્થાન
તે મુજબ સુધાર્યું છે.}
$x$ એ $!A$માં મુક્ત
ન હોવાથી
$\Subst{!A}{x}{c} \Subst{}{y}{x} = \Subst{!A}{y}{c}$;
તેથી
\[
\Proves \lforall[x][\Subst{!A}{x}{c}] \lif \Subst{!A}{y}{c}, 
\] 
અને નિગમન પ્રમેયથી
$\lforall[x][\Subst{!A}{x}{c}] \Proves
\Subst{!A}{y}{c}$ મળે.
$y$ એ $!A$માં મુક્ત
ન હોવાથી
$\lforall[x][\Subst{!A}{x}{c}]$માં
પણ મુક્ત નથી; તેથી
સર્વત્રીકરણથી
$\lforall[x][\Subst{!A}{x}{c}] \Proves \lforall[y][\Subst{!A}{y}{c}]$,
અને ફરી નિગમન
પ્રમેયથી $\Proves
\lforall[x][\Subst{!A}{x}{c}] \lif \lforall[y][\Subst{!A}{y}{c}]$.
$\Proves \lforall[y][\Subst{!A}{y}{c}] \lif \lforall[x]
[\Subst{!A}{x}{c}]$ની
સાબિતી સંપૂર્ણ
સમમિત છે; તેથી
ટોટોલોજી વડે
નિષ્કર્ષ મળે છે.
\end{proof}

\begin{thm}[અચળો પર પ્રબળ સર્વત્રીકરણ]
\ollabel{thm:StrongGen} 
જો $\Gamma \Proves !A$, !!{constant}
$c$ એ $\Gamma$માં આવતું ન હોય,
અને $x$ એ $!A$માં મુક્ત
ન હોય છતાં $!A$માં
$c$ની જગ્યાએ મૂકવા
!!{free for} હોય, તો
$\Gamma \Proves \lforall[x][\Subst{!A}{x}{c}]$.
\end{thm}

\begin{proof}
કારણ કે $\Gamma \Proves !A$
અને $c$ એ $\Gamma$માં
નથી, નબળા સર્વત્રીકરણથી
$!A$માં ન આવતું
!!a{variable} $y$ મળે કે
$\Gamma \Proves \lforall[y][\Subst{!A}{y}{c}]$.
$y$ એ $!A$માં બિલકુલ
ન આવતું હોવાથી તે
$!A$માં મુક્ત નથી
અને $!A$માં $c$ની જગ્યાએ
મૂકવા !!{free for} પણ છે.
વળી, પરિકલ્પના મુજબ
$x$ એ $!A$માં
મુક્ત નથી અને
$!A$માં $c$ની જગ્યાએ
મૂકવા !!{free for} છે.
એટલે બંધ !!{variable}
બદલવાની શરતો પૂરી થાય છે:
$\Proves \lforall[x][\Subst{!A}{x}{c}] \equiv
\lforall[y][\Subst{!A}{y}{c}]$;
તેથી $\Gamma \Proves
\lforall[x][\Subst{!A}{x}{c}]$.
\end{proof}

\begin{thm} 
\ollabel{thm:provability-quantifiers}
\begin{tagenumerate}{prvEx,prvAll} 
\tagitem{prvEx}{જો $\Gamma \Proves !A(t)$, તો $\Gamma \Proves
  \lexists[x][!A(x)]$.}{} 
\tagitem{prvAll}{જો $\Gamma \Proves \lforall[x][!A(x)]$, તો $\Gamma
  \Proves !A(t)$.}{}
\end{tagenumerate} 
\end{thm}

\begin{proof} 
\begin{tagenumerate}{prvEx,prvAll} 
%%Need axioms for exists
\tagitem{prvEx}{કસરત.}{}
\tagitem{prvAll} {માનો કે $\Gamma_0 \Proves \lforall[x][!A(x)]$.

\begin{derivation} 
1. & $\Gamma_0 \Proves \lforall[x][!A(x)]$ & ધારણા \\
2. & $\Gamma_0 \Proves \lforall[x][!A(x)] \lif \Subst{!A}{t}{x}$ & Ax1 \\ 
3. & $\Gamma_0 \Proves \Subst{!A}{t}{x}$ & MP 1, 2 \\ 
\end{derivation}}{}

\end{tagenumerate} 
\end{proof}

\end{document}
